% sections/08_limitations.tex

\section{Limitations}
\label{sec:limitations}

\subsection{Layer 2 Not Yet Implemented}

The evaluation framework described in this paper is a Layer-1 harness: it
imports tool functions directly, seeds the random state, and measures
correctness deterministically without a live language model in the loop.
A Layer-2 evaluation---where a real LLM agent calls the MCP server over the
protocol and its outputs are scored---has not been implemented.
Consequently, there are no live-agent result numbers in this paper.
All result rows marked \texttt{\textbackslash TODO} in the tables of
Section~\ref{sec:results} require the Layer-2 harness to be built and run.
The Layer-1 results establish that the deterministic core is correct; they
do not establish that a given LLM will interact with the server correctly
in end-to-end use.

\subsection{Small Dataset: S4 FTIR Purity}

Scenario S4 (FTIR purity classification) uses approximately 20 spectra
assigned to 7 distinct groups.
At this sample count, cross-validation fold assignments are highly sensitive
to group ordering, per-class counts are below typical minimum thresholds,
and covariance-based models (PCA-LDA) encounter singular matrices.
The case study in Section~\ref{sec:casestudy} documents exactly these
failure modes.
Results from S4 are qualitatively informative about HITL behavior under
data scarcity, but they should not be treated as representative of FTIR
classification performance at scale.

\subsection{Single Instrument and Single FTIR Source}

NIR measurements in scenarios S1--S3 were collected with a single
trinamiX handheld instrument; the FTIR spectra in S4 come from a single
laboratory source.
Instrument-specific effects (detector response, scan averaging, spectral
resolution) are not separated from chemical variance in any of the scenarios.
The modality registry enforces axis-unit and value-domain constraints for
the instruments represented, but these constraints have not been validated
against spectra from other instruments of the same modality type.
External validity---whether a model trained on trinamiX NIR data will
generalize to Bruker or Perkin-Elmer NIR data, for example---is outside the
scope of this paper.

\subsection{Group-Label Plumbing Requires Verification}

The \texttt{run\_analysis} tool accepts a \texttt{group\_column} argument
that should propagate through the entire pipeline to ensure grouped
cross-validation is applied when replicate structure is present.
End-to-end plumbing of this label---from user-provided data through
preprocessing, CV splitting, and leakage detection---has been designed but
not fully verified across all code paths.
A regression could silently revert to naive KFold if the group column is
dropped during data transformation.
This is a known gap targeted for verification before Layer-2 evaluation.

\subsection{Hallucination Probes Are Hand-Authored and Limited}

The hallucination guard fixtures used in Layer-1 evaluation were authored
by hand and cover 6 distinct probe scenarios.
The probes were designed to represent common LLM failure modes (metric
fabrication, causal overclaiming, silent fallback acceptance), but they are
not exhaustive.
New LLM behaviors could produce hallucinations outside the coverage of
the existing fixtures without being detected by the current harness.
Expanding the probe library to cover additional failure modes, and automating
probe generation, is deferred to future work.

\subsection{Expert Baselines Are Minimal}

Expert baselines (what a human chemometrics practitioner would produce
without the tool system) consist of a single model with default
hyperparameters and no tuning.
No domain-expert cross-validation of the baseline choices was performed,
and inter-rater reliability on plan quality scores has not been assessed.
As a result, comparisons between the agent output and the expert baseline
should be interpreted cautiously.

\subsection{Plan Quality Score Is Automated Only}

Plan quality is scored by a 4-boolean automated rubric (presence of
preprocessing choice, rationale, CV strategy, and metric justification).
This rubric captures structural completeness but not scientific soundness.
A human expert reviewing the same plans on a 1--5 scale might assign
meaningfully different rankings.
Human scoring and inter-rater reliability assessment are deferred to the
Layer-2 evaluation.
